\chapter{Git on your own}
\label{ch:git}

\chaptergoals{1 of 6 in Part~I: the first of three git chapters.}
{a terminal, and \cmd{git} (\cmd{git -{}-version} prints 2.23 or later).
Nothing else: no GitHub account, no SHiP software.}
{keep the history of your own work: make commits, look back at them, undo
mistakes, and keep data out of the repository.}

Every SHiP package lives in a git repository on GitHub, under
\href{https://github.com/ShipSoft}{\texttt{github.com/ShipSoft}}, and every
change to it arrives as a pull request. Part~I builds up to that in small
steps. This chapter is git on your own machine only; \cref{ch:branches} adds
branches, \cref{ch:remotes} other people and GitHub, \cref{ch:pixi,ch:ci} the
environment and the automatic checks around it, and \cref{ch:team} puts
everything together.

\section{What git keeps}

A git repository is a history of \emph{snapshots} of a directory. Each
snapshot is a \emph{commit}: the complete state of every tracked file, an
author, a date, a message, and a pointer to the commit it was made from (its
\emph{parent}). A commit is named by its hash, such as \texttt{c8dfe70}, which
is computed from all of this, so a commit can never change: ``changing'' one
means making a new commit with a new hash.

You work in three places:

\begin{center}
\begin{tikzpicture}[
  area/.style={draw=shipnavy, fill=shipcream, rounded corners=2pt, minimum width=32mm,
    minimum height=11mm, align=center, font=\sffamily\small},
  move/.style={-{Stealth[length=2mm]}, line width=0.7pt, color=shipnavy},
  back/.style={move, dashed},
  lab/.style={font=\ttfamily\scriptsize, inner sep=1.5pt, fill=white}]
  \node[area] (wt) {working tree\\\scriptsize the files you edit};
  \node[area, right=22mm of wt] (ix) {staging area\\\scriptsize the next commit};
  \node[area, right=22mm of ix] (repo) {repository\\\scriptsize all commits};
  \draw[move] (wt) -- node[lab, above] {git add} (ix);
  \draw[move] (ix) -- node[lab, above] {git commit} (repo);
  \draw[back] (ix.south) to[bend left=20] node[lab, pos=0.5] {git restore -{}-staged} (wt.south);
  \draw[back] (repo.south) to[bend left=32] node[lab, pos=0.5] {git restore -{}-source} ([xshift=-6mm]wt.south);
\end{tikzpicture}
\end{center}

\texttt{git add} chooses what goes into the next commit; \texttt{git commit}
records it. Because you choose, one commit can hold one logical change even if
you edited several things.

\section{Editing files}
\label{sec:editor}

The commands in this booklet open files with \cmd{vim}, which is on every
Linux machine, including lxplus. Any editor will do (nano, VS Code, Emacs,
\ldots): replace \cmd{vim} by yours. What you need of vim to start:

\begin{center}\small
\begin{tabular}{@{}l l@{}}
\toprule
\cmd{vim notes.md} & open the file (vim creates it when you save) \\
\cmd{i}            & insert mode: now type the text \\
\cmd{Esc}          & back to normal mode, where keys are commands \\
\cmd{:w}, \cmd{:q} & save; quit \\
\cmd{:wq}          & save and quit \\
\cmd{:q!}          & quit and throw the changes away \\
\cmd{u}, \cmd{dd}, \cmd{/word} & undo; delete the line; search \\
\bottomrule
\end{tabular}
\end{center}

\begin{pitfall}
Stuck in vim? Press \cmd{Esc}, type \cmd{:q!} and press Enter. Then run
\cmd{vimtutor} once, in a spare half hour: it teaches the rest in 30 minutes.
\end{pitfall}

git opens the same editor whenever it needs text from you (a commit message
without \cmd{-m}, a rebase to continue). The first exercise tells it which
one.

\section{First steps}

\begin{exercise}{Your first repository}
\label{ex:gitfirst}
Tell git who you are (once per machine), then make a repository:
\begin{shell}
git config --global user.name "Your Name"
git config --global user.email "your.name@cern.ch"
git config --global init.defaultBranch main
git config --global pull.rebase true
git config --global core.editor vim
mkdir git-playground && cd git-playground
git init
vim hello.py                     # write print('hello'), then :wq
git status                       # hello.py is "untracked"
git add hello.py
git commit -m "feat: say hello"
git log --oneline
\end{shell}
\end{exercise}

\begin{note}
Use the e-mail address of your GitHub account, so that GitHub can attach your
commits to you later. \cmd{pull.rebase true} only matters from
\cref{ch:remotes} on; \cref{sec:mergerebase} explains why SHiP wants it.
\end{note}

\begin{exercise}{Stage, look, undo}
Open \file{hello.py} in your editor again and change it to print
\texttt{'hello, SHiP'}.
\begin{enumerate}
  \item \cmd{git status} and \cmd{git diff}: what does git say changed?
  \item \cmd{git add hello.py}, then \cmd{git diff} and
        \cmd{git diff -{}-staged}. Why is the first one empty now?
  \item Unstage with \cmd{git restore -{}-staged hello.py}, then throw the
        change away with \cmd{git restore hello.py}. Check with
        \cmd{cat hello.py}.
\end{enumerate}
\end{exercise}

\begin{pitfall}
\cmd{git restore <file>} discards your uncommitted changes to that file for
good: git never saw them, so it cannot give them back. What has been
committed, on the other hand, is very hard to lose.
\end{pitfall}

\section{Good commits}

\paragraph{One change per commit.} A commit should do one thing that you can
describe in one line: easier to review, to understand later, and to undo.

\paragraph{Commit messages.} The SHiP repositories follow
\href{https://www.conventionalcommits.org}{Conventional Commits}: a type, an
optional scope, and a short summary in the imperative.
\begin{center}\small
\begin{tabularx}{\linewidth}{>{\ttfamily}l L}
\toprule
feat(task-07): fit with a Poisson likelihood & a new feature \\
fix(field): place the map at the magnet centre & a bug fix \\
docs: explain the sensitive-volume order & documentation only \\
refactor:\ test:\ style:\ chore: & restructuring, tests, formatting, housekeeping \\
\bottomrule
\end{tabularx}
\end{center}
The release changelogs are generated from these messages, and in the ShipSoft
repositories an automatic check rejects pull requests whose commits do not
follow the convention (\cref{ch:ci}).

\begin{exercise}{Read the history}
Make two more commits in \file{git-playground}: change the message printed by
\file{hello.py}, and add a second file. Then:
\begin{shell}
git log --oneline              # one line per commit, newest first
git show HEAD                  # the latest commit: message and change
git log -p hello.py            # every change ever made to one file
git diff HEAD~2 HEAD           # everything between two commits ago and now
\end{shell}
\texttt{HEAD} is the commit you are on; \texttt{HEAD\textasciitilde1} its
parent, \texttt{HEAD\textasciitilde2} its grandparent. Get an old version of a
file back with \cmd{git restore -{}-source HEAD\textasciitilde2 hello.py},
then commit it, or restore it again from \texttt{HEAD}.
\end{exercise}

\section{What not to commit}

Data, plots and build products do not belong in git: they are large, they
change with every run, and they can be made again from the code. List them
in a file called \file{.gitignore}, and git stops offering them to you.

\begin{exercise}{Ignore the outputs}
\begin{shell}
touch data.root plot.pdf
git status                     # both are offered as "untracked"
vim .gitignore                 # two lines: *.root and *.pdf
git status                     # now only .gitignore is
git add .gitignore && git commit -m "chore: ignore data and plots"
\end{shell}
\end{exercise}

\begin{levelcheck}
\begin{itemize}
  \can make a repository, stage exactly the changes you want, and commit them
       with a Conventional-Commit message;
  \can see what changed (\cmd{status}, \cmd{diff}, \cmd{diff -{}-staged}) and
       read the history (\cmd{log}, \cmd{show});
  \can undo an uncommitted change, and get an old version of a file back;
  \can keep data and outputs out of the repository with \file{.gitignore}.
\end{itemize}
\end{levelcheck}
